\section{Sintaks Dasar HTML: Tag, Elemen, dan Atribut}

HTML dibangun dari tiga komponen fundamental: tag, elemen, dan atribut. Pemahaman mendalam tentang ketiga komponen ini menjadi dasar bagi setiap pengembang web untuk menulis markup yang valid, terstruktur, dan mudah dipelihara. \textbf{Tag} adalah penanda yang ditulis dalam tanda kurung sudut (angle brackets) seperti \code{<p>}, \code{<h1>}, atau \code{<div>}. Sebagian besar elemen HTML terdiri dari tag pembuka dan tag penutup yang membungkus konten; tag pembuka menandakan awal elemen, sementara tag penutup yang ditandai dengan garis miring (misalnya \code{</p>}) menandakan akhir elemen \cite{mdnHtmlIntro}.

\textbf{Elemen} adalah keseluruhan unit yang terdiri dari tag pembuka, konten di dalamnya, dan tag penutup. Sebagai contoh, dalam \code{<p>Ini paragraf.</p>}, keseluruhan unit dari awal \code{<p>} hingga akhir \code{</p>} disebut sebagai elemen paragraf. Terdapat juga kategori khusus yang disebut \textit{void elements} atau elemen kosong yang tidak memiliki konten dan oleh karena itu tidak memerlukan tag penutup. Contoh void elements meliputi \code{<img>}, \code{<br>}, \code{<hr>}, \code{<input>}, dan \code{<meta>}. Pada HTML5, void elements dapat ditulis dengan atau tanpa garis miring di akhir (misalnya \code{<img />} atau \code{<img>}), namun XHTML memerlukan garis miring untuk validitas \cite{whatwgHtml}.

\textbf{Atribut} menyediakan informasi tambahan tentang elemen dan ditulis di dalam tag pembuka dengan format \code{nama="nilai"}. Atribut memungkinkan kita mengkustomisasi perilaku atau memberikan metadata pada elemen. Beberapa atribut bersifat global dan dapat digunakan pada semua elemen HTML, seperti \code{id} (identifier unik), \code{class} (untuk klasifikasi styling), \code{style} (CSS inline), \code{title} (tooltip), \code{lang} (bahasa konten), dan \code{data-*} (atribut data kustom). Atribut lain bersifat spesifik untuk elemen tertentu, seperti \code{href} untuk link, \code{src} untuk gambar, dan \code{type} untuk input.

\begin{table}[htbp]
\centering
\begin{tabular}{|P{3cm}|P{3cm}|P{6cm}|}
\hline
\textbf{Konsep} & \textbf{Contoh} & \textbf{Penjelasan} \\
\hline
Tag & \code{<p>}, \code{</p>} & Penanda markup dalam tanda kurung sudut \\
\hline
Elemen & \code{<p>Teks</p>} & Tag pembuka + konten + tag penutup \\
\hline
Atribut & \code{<p class="intro">} & Informasi tambahan di tag pembuka \\
\hline
Void Element & \code{<img />}, \code{<br />} & Elemen tanpa konten dan tag penutup \\
\hline
\end{tabular}
\caption{Perbandingan Tag, Elemen, dan Atribut dalam HTML}
\label{tab:tag-element-attr}
\end{table}

Aturan \textit{nesting} (penyusunan bersarang) adalah aspek penting dalam struktur HTML yang menentukan bagaimana elemen dapat berada di dalam elemen lain. HTML mengikuti model hierarki pohon (tree) di mana elemen dapat memiliki elemen anak (child elements). 
\begin{sloppypar}
Aturan nesting mengharuskan tag penutup ditulis dalam urutan terbalik dari tag pembuka - prinsip yang dikenal sebagai \textit{LIFO} (Last In, First Out). Misalnya, \code{<p><strong>Teks</strong></p>} adalah valid, sementara \code{<p><strong>Teks</p></strong>} adalah invalid karena tag ditutup dalam urutan yang salah. Nesting yang benar memastikan struktur DOM (Document Object Model) terbentuk dengan baik, yang sangat penting untuk CSS styling dan JavaScript manipulation \cite{mdnHtmlIntro}.
\end{sloppypar}

\begin{webcode}[language=HTML, caption={Contoh Tag, Elemen, dan Atribut}]
<!-- Elemen paragraf dengan atribut class -->
<p class="intro" id="paragraf-utama">
  Ini adalah paragraf dengan <strong>teks tebal</strong>.
</p>

<!-- Link dengan multiple atribut -->
<a href="https://example.com" 
   target="_blank" 
   rel="noopener noreferrer"
   title="Kunjungi situs contoh">
  Link ke Example.com
</a>

<!-- Gambar (void element) dengan atribut lengkap -->
<img src="foto.jpg" 
     alt="Deskripsi foto" 
     width="400" 
     height="300" />
\end{webcode}

\begin{catatan}
\textbf{Aturan Nesting yang Benar:}
\begin{itemize}[leftmargin=*]
  \item Selalu tutup tag dalam urutan terbalik dari pembukaan (LIFO - Last In, First Out).
  \item Elemen blok (block-level) seperti \code{<div>} dapat berisi elemen inline dan blok lainnya.
  \item Elemen inline seperti \code{<span>} atau \code{<a>} tidak boleh berisi elemen blok.
  \item Gunakan indentasi konsisten (2 atau 4 spasi) untuk memudahkan pembacaan struktur bersarang.
\end{itemize}
\end{catatan}

